\section{Analytic prices from real coupling estimates}
\label{sec:real-coupling-interpolation}
The parameter estimates used for the background history have a useful
feature: they are obtained from real configurations. The local weights
contain exponentials whose exponents need not be bounded, and bounding
their absolute values before using the real estimates would lose the
required smallness. We prove the interpolation argument that avoids
this loss. The real local estimates to which it is applied retain their
full derivative and support weights.

\subsection{An algebra of parameter series}
For a real power series $f(t)=\sum_{r\ge0}f_rt^r$ and $R\ge1$, set
\[
                         |f|_R=\sum_{r\ge0}|f_r|R^r.
\]
Only series with finite displayed norm will be used. Their products
are given by the absolutely convergent Cauchy product, and
$|fg|_R\le |f|_R|g|_R$. Consequently
$|\exp f|_R\le\exp(|f|_R)$. If the constant coefficient $f_0$
is real, it is better to keep its sign:
\begin{equation}\label{eq:real-series-exponential}
                 |\exp f|_R\le e^{f_0}\exp(|f-f_0|_R).
\end{equation}
These assertions follow by summing the nonnegative absolute
coefficient bounds for $\sum_{m\ge0}f^m/m!$; absolute convergence
also identifies that series with the exponential on $[-R,R]$.
For a finite real test $\lambda=\sum_ha_h\delta_{t_h}$,
$|t_h|\le1$, put
\[
             M_R(\lambda)=\sup_{r\ge0}
                          \left|\sum_h a_h(t_h/R)^r\right|.
\]
Termwise summation then gives
\begin{equation}\label{eq:real-series-test}
                    \left|\sum_h a_h f(t_h)\right|
                                      \le |f|_R M_R(\lambda).
\end{equation}
Thus an absolute coefficient estimate supplies every finite-test
estimate simultaneously.

Fix a real normed space $E$, a point $x\in E$, and an integer $j\ge0$.
For a $C^j$ function $V$ near $x$, write
$D_u^iV(x)=\left.\frac{d^i}{ds^i}V(x+su)\right|_{s=0}$.
A diagonal jet bound with amplitude $C$ and derivative scale $d$
means $|D_u^iV(x)|\le C d^i$ for $0\le i\le j$ and
$\|u\|\le1$. When $i=0$ we use $d^0=1$, including when $d=0$.
The parameter estimates below hold for each such derivative, with
one common amplitude. Conversion to full multilinear derivative
norms uses the fixed finite-order polarization constants already
included in the local derivative prices.

\subsection{The larger real nodes}
Let $\ell(t)=a+bt$, and suppose
\begin{equation}\label{eq:real-affine-band}
 R\ge1,\qquad 1\le a-6R|b|,\qquad
 a+6R|b|\le1+\delta,\qquad \delta\le1.
\end{equation}
In particular $0\le\delta\le1$, and
$1\le\ell(t)\le1+\delta$ whenever $|t|\le6R$.
Set $p(t)=1-\ell(t)^2$. For every real $q$ define
\[
 A(q)=(1-a^2)q+2a|bq|R+b^2|q|R^2.
\]
By \eqref{eq:real-series-exponential},
\begin{equation}\label{eq:real-exponent-price}
             |\ell|_R\le1+\delta,\qquad
             |\exp(pq)|_R\le e^{A(q)}.
\end{equation}
The first inequality uses $a\ge1$ and
$a+R|b|\le a+6R|b|$. No estimate on $|q|$ is required.

Choose $\sigma=-1$ if $qb\ge0$, and $\sigma=1$ otherwise. For
$r=0,\ldots,2j$, put
\[
 u_r=3+\frac{r}{2j+1},\qquad \xi_r=\sigma u_r,
                   \qquad t_r=R\xi_r.
\]
These are distinct real points in $[-6R,6R]$. At every one of them,
\begin{equation}\label{eq:real-node-domination}
                             A(q)\le p(t_r)q.
\end{equation}
Here is the sign-sensitive estimate. Put $d_*=R|b|$; then
$a\ge1+6d_*$ and $qbR\sigma=-|q|d_*$. For any
$3\le u\le6$, we have
\[
 2ad_*+d_*^2\le3ad_*\le ad_*u
                              \le2ad_*u-d_*^2u^2.
\]
The last inequality is $d_*u\le a$, multiplied by $d_*u$;
the other two follow from $d_*\le a$ and $u\ge3$.
Expanding $p(R\sigma u)q$, and using $q\le|q|$, now gives
\[
\begin{split}
 p(R\sigma u)q
  &=(1-a^2)q+2a|q|d_*u-qd_*^2u^2\\
  &\ge(1-a^2)q+|q|(2ad_*u-d_*^2u^2)\ge A(q).
\end{split}
\]
This proves \eqref{eq:real-node-domination} for either sign of $q$.

\subsection{Interpolating the local weight with its derivatives}
Let $L,\Gamma:E\to\mathbb R$ be $C^j$ near $x$, and let $n$ be a
nonnegative integer. Consider
\begin{equation}\label{eq:real-local-weight}
             J(t,y)=L(y)\ell(t)^n e^{p(t)\Gamma(y)}.
\end{equation}
Choose the preceding nodes with $q=\Gamma(x)$. Suppose that at those
$2j+1$ nodes the real local estimates give
\begin{equation}\label{eq:real-node-jet}
               |D_u^iJ(t_r,x)|\le C d^i
 \quad(0\le i\le j,\ \|u\|\le1),
                      \qquad C,d\ge0.
\end{equation}
There is a constant $K_j$, depending only on $j$, such that each
function $t\mapsto D_u^iJ(t,x)$ has an absolutely convergent series
with
\begin{equation}\label{eq:real-local-series-price}
             |D_u^iJ(\cdot,x)|_R
                      \le K_j(1+\delta)^n C d^i.
\end{equation}
In particular this controls all finite tests by
\eqref{eq:real-series-test}. The constant is independent of
$a,b,R,n,C,d$ and of the value of $\Gamma(x)$.

We prove the assertion, including the reason that only finitely many
real nodes suffice. Differentiating along $x+su$ and factoring out
the exponential gives
\begin{equation}\label{eq:real-polynomial-jet}
 D_u^iJ(t,x)=P_{i,u}(t)\ell(t)^n e^{p(t)\Gamma(x)},
                           \qquad \deg P_{i,u}\le2i.
\end{equation}
Indeed, for the exponential alone its $i$th derivative is the
exponential times a sum, over partitions of $\{1,\ldots,i\}$,
of products $\prod_B p(t)D_u^{|B|}\Gamma(x)$. Each product has
at most $i$ factors $p(t)$, each of degree two. Leibniz's rule
with $L$ gives \eqref{eq:real-polynomial-jet}, also for $i=0$.
No division by $L$ or by a derivative of $\Gamma$ occurs.

Let $\mathcal L_r(z)$ be the Lagrange basis polynomials for the
nodes $\xi_0,\ldots,\xi_{2j}$. Polynomial interpolation gives
\[
          P_{i,u}(t)=\sum_{r=0}^{2j}P_{i,u}(t_r)
                                      \mathcal L_r(t/R).
\]
Define $K_j$ as the larger, for the two possible signs $\sigma$,
of $\sum_r|\mathcal L_r|_1$. This is finite, since all the nodes
are distinct. Also $|\mathcal L_r(\cdot/R)|_R
=|\mathcal L_r|_1$. By \eqref{eq:real-node-jet},
\eqref{eq:real-polynomial-jet}, and $\ell(t_r)^n\ge1$,
\[
        |P_{i,u}(t_r)|e^{p(t_r)q}\le C d^i.
\]
Multiplying each interpolation term by
$\ell(t)^n e^{p(t)q}$ and using
\eqref{eq:real-exponent-price} and
\eqref{eq:real-node-domination}, its coefficient norm is at most
\[
 |P_{i,u}(t_r)|(1+\delta)^n e^{A(q)}
                    |\mathcal L_r|_1
       \le (1+\delta)^n C d^i|\mathcal L_r|_1.
\]
Summing proves \eqref{eq:real-local-series-price}. Notice that the
node sign may depend on $\Gamma(x)$. Both sets of nodes belong to
the same allowed real interval, and $K_j$ controls either choice;
there is no need to differentiate the choice of sign.

\subsection{The Gaussian density and the local support count}
For a fixed Gaussian coordinate $z$, let $Q(z)\ge0$ be its quadratic
energy in the reference density. Under the rescaling
$z\mapsto\ell(t)z$ in $D$ real dimensions, the density ratio,
including the Jacobian, is
\[
                 g(t,z)=\ell(t)^D e^{(1-\ell(t)^2)Q(z)}.
\]
The same argument, now with $q=Q(z)\ge0$, yields
\begin{equation}\label{eq:real-gaussian-series-price}
                              |g(\cdot,z)|_R\le(1+\delta)^D.
\end{equation}
Indeed \eqref{eq:real-node-domination} gives
$A(q)\le p(t_r)q\le0$, since $\ell(t_r)\ge1$.
The bound is uniform in $z$. Thus it can be multiplied into a
local coefficient majorant before integration against the fixed
reference Gaussian measure.

In the local application, the exponent $n$ in
\eqref{eq:real-local-weight} counts the seed coordinates actually
present in the local weight, each with its color dimension. The
Gaussian dimension $D$ is likewise that of the relevant marginal.
These counts are determined by the local support; replacing them
by the dimension of the entire torus would destroy a bound uniform
in volume. The exact local change-of-density identity is used
before applying \eqref{eq:real-gaussian-series-price}. The remaining
support-dependent powers $(1+\delta)^{n+D}$ are included in the
cover-price reserve.

\subsection{Keeping the small amplitude of ordinary activities}
There is a second, complementary estimate. Let
$f:\mathbb R\times E\to\mathbb R$ be real-valued, and suppose that
$t\mapsto f(t,y)$ is a polynomial of degree at most $k$ on
$[-1,1]$ for every $y$, where $k$ is fixed. Choose distinct fixed
nodes $s_0,\ldots,s_k\in[-1,1]$. Suppose that each $f(s_r,\cdot)$
is $C^j$ near $x$ and that
\begin{equation}\label{eq:real-small-activity}
 |D_u^i(e^{f(s_r,\cdot)}-1)(x)|\le v d^i,
                  \qquad 0\le v\le\tfrac12.
\end{equation}
There are constants $A,B\ge1$, depending only on $j,k$ and the
nodes, such that, whenever $BR^kv\le1$,
\begin{equation}\label{eq:real-small-series}
        |D_u^i(e^{f(\cdot,\cdot)}-1)(x)|_R
                                      \le AR^k v d^i.
\end{equation}
In this formula the derivative is in the second variable and the
series norm in the first.

To check the estimate, write $h=e^{f(s_r,\cdot)}-1$ at a fixed
node. At $x$ we have $|h(x)|\le v\le1/2$ and
$f(s_r,\cdot)=\log(1+h)$ near $x$. The zeroth derivative is
bounded by $2v$, by integrating $1/(1+s)$ between $0$ and $h(x)$.
For a positive-order derivative, the chain rule is a finite sum
over partitions, with one derivative of $\log(1+s)$ and products
of derivatives of $h$. The former is bounded on $[-1/2,1/2]$;
a term with $r\ge1$ factors of $h$ is bounded by a constant
times $v^r d^i\le v d^i$. Hence
$|D_u^if(s_r,\cdot)(x)|\le C_jv d^i$ at every node.

Polynomial interpolation in $t$ is valid for every $y$ and so
also after the indicated derivatives. The coefficient norm of
any fixed degree-$k$ Lagrange basis polynomial at radius $R\ge1$
is bounded by $R^k$ times its norm at radius one. It follows that
\[
                  |D_u^if(\cdot,x)|_R\le C_*R^k v d^i.
\]
Choose $B\ge C_*$. Put $\varepsilon=C_*R^kv\le1$. In the
Banach algebra of parameter series, the zeroth-order estimate
$|e^f-1|_R\le e^\varepsilon-1\le e\varepsilon$ follows from
the exponential series. For $i>0$, the same partition formula
for derivatives of $e^f$ gives a factor $|e^f|_R\le e$ and
products bounded by $\varepsilon^r d^i\le\varepsilon d^i$.
The finite sum gives \eqref{eq:real-small-series}, with $A$
large enough for all $i\le j$. This proof also covers $d=0$:
all positive derivatives then vanish.

The incoming ordinary activity in the coupling family has degree
at most two in this real parameter. Thus the radius cost in this
last argument is $R^2$, with constants fixed by the derivative
order. The local stock weight uses the preceding real-node
argument instead, because its exponent need not be small.
Together with the local real jet estimates, the fixed Gaussian
change of density, and the summable support prices, these two
arguments provide the parameter estimates for the finite forest
integrals. Summation over connected supports and the subsequent
local projections still require their stated quantitative bounds;
the interpolation proved here does not replace those bounds.
